\section{故障诊断实验与结果分析}
\label{sec:experiments}

\subsection{引言}

前述章节完成了汽轮机系统健康状态建模、故障知识图谱构建以及KG-GCN故障诊断模型设计。本章选取超高压缸（VHP）作为代表性设备，对本文提出的汽轮机系统KG-GCN故障诊断方法开展算例验证。考虑到当前机组尚缺少能够覆盖多类典型故障的完整现场故障样本，本文在真实健康运行残差基础上，依据独立故障机理构造半物理故障场景，用于验证不同诊断模型对典型汽轮机故障模式的识别能力。实验数据不作为现场真实故障记录使用，相关结果用于说明所提方法在机理一致构造场景下的可行性。

\subsection{研究对象及故障数据集}
\label{subsec:fault_dataset_v4}

本文研究对象为某燃煤机组汽轮机热力系统。第三章构建的系统级知识图谱覆盖主蒸汽、VHP、一次再热、HP、二次再热、IP、LP、抽汽回热以及凝汽器冷端等主要设备和参数关系。为验证所提KG-GCN方法，本文选取VHP作为典型设备开展代表性算例。VHP一方面直接处于主蒸汽膨胀通流路径中，另一方面通过一级抽汽支路与1号高压加热器相连，其排汽又构成一次再热系统入口边界，因此在单一设备范围内即可形成VHP本体、抽汽支路和再热接口等多个具有明确物理位置的状态区域。

第二章VHP健康状态模型最终保留24个状态参数作为诊断节点，包括3个VHP叶片级压力、排汽压力和温度、一级抽汽口压力与温度、一级抽汽逆止门前后温度、1号高压加热器入口抽汽压力与温度、一次低温再热接口温度以及VHP缸壁和内部蒸汽温度。上述参数均具有固定DCS点号，并通过第三章设备--参数关系映射至VHP任务参数子图。

结合汽轮机通流故障机理、汽封泄漏机理、抽汽支路结构以及当前24个DCS状态参数的可观测范围，构造正常状态和三类典型故障场景，如表~\ref{tab:fault_classes_yangstyle}所示。其中，具体测点响应幅值和时间演化形式仅用于半物理故障样本构造，不反向写入KG邻接矩阵，以避免样本生成知识与分类拓扑之间形成循环论证。

\begin{table}[htbp]
    \centering
    \caption{VHP故障诊断样本类别}
    \label{tab:fault_classes_yangstyle}
    \begin{tabular}{p{1.0cm}p{4.2cm}p{5.2cm}p{3.3cm}}
        \toprule
        类别 & 故障场景 & 主要状态特征 & 主要作用区域 \\
        \midrule
        F0 & 正常状态 & 真实健康运行残差，不叠加故障响应 & 全部诊断节点 \\
        F1 & VHP通流部件侵蚀 & 通流有效面积及级组压力分布发生改变 & VHP本体、排汽及再热接口 \\
        F2 & 一级抽汽支路泄漏 & 抽汽支路蒸汽损失，下游抽汽状态衰减 & 一级抽汽及1号高加入口 \\
        F3 & VHP汽封泄漏 & 无效蒸汽泄漏引起本体及排汽热力状态改变 & VHP本体、排汽及相关接口 \\
        \bottomrule
    \end{tabular}
\end{table}

健康背景数据取自2024年6月独立运行数据。利用第二章GRU健康状态模型计算24个状态参数的实测值与健康预测值之差，并以此形成真实健康残差背景。对于第$k$类故障，第$n$个半物理故障残差窗口表示为
\begin{equation}
R_{k,n}^{F}=R_n^{H}+a_{k,n}S_k\odot G_{k,n}+\epsilon_{k,n},
\end{equation}
式中：$R_n^{H}$为真实健康残差窗口；$S_k$为由独立故障机理确定的空间响应模式；$a_{k,n}$为故障严重程度系数；$G_{k,n}$为故障随时间演化的动态函数；$\epsilon_{k,n}$为小幅随机扰动。不同样本随机设置故障发生时刻，并采用阶跃、斜坡或Sigmoid等不同演化形式，避免模型仅记忆单一固定故障模板。主实验中故障强度在训练健康残差标准差的$2.0\sim3.0\sigma$范围内随机变化，该范围表示可明确观测的异常残差强度，并不对应设备实际损伤百分比。

为避免相邻窗口造成数据泄漏，首先按照自然日对真实健康背景进行互斥划分，再分别在训练、验证和测试日期内部构造故障窗口。每连续5个自然日中，第1--3日作为训练集，第4日作为验证集，第5日作为测试集，且任一60~min窗口不得跨越自然日边界。按照上述规则，训练集、验证集和测试集分别包含1712、556和560个样本，各故障类别在不同数据子集中保持均衡，具体如表~\ref{tab:fault_dataset_yangstyle}所示。

\begin{table}[htbp]
    \centering
    \caption{故障诊断数据集样本分布}
    \label{tab:fault_dataset_yangstyle}
    \begin{tabular}{lrrrr}
        \toprule
        数据集 & F0 & F1 & F2 & F3 \\
        \midrule
        训练集 & 428 & 428 & 428 & 428 \\
        验证集 & 139 & 139 & 139 & 139 \\
        测试集 & 140 & 140 & 140 & 140 \\
        \bottomrule
    \end{tabular}
\end{table}

\subsection{实验设置与评估指标}
\label{subsec:experiment_setting_yangstyle}

为验证本文所提KG-GCN模型的故障诊断性能，选取CNN和AE作为对比算法。CNN将多参数残差窗口作为二维输入，通过局部卷积核提取参数--时间特征；AE首先通过无监督重构学习低维潜在特征，再利用分类层完成故障识别；KG-GCN则将状态残差加载至知识图谱对应参数节点，并依据第三章构建的参数关联路径进行图卷积特征传播。三种模型使用完全相同的训练、验证和测试数据划分。

KG-GCN输入节点数为24，残差窗口长度为60，图卷积编码器包含2层GCN，每个节点的隐藏特征维数为32。无标签预训练阶段随机遮蔽15\%的残差位置，通过重构被遮蔽残差训练图卷积编码器；完成预训练后固定GCN编码器参数，并训练全连接分类层。优化器采用AdamW，权重衰减为$1\times10^{-4}$。主要实验参数见表~\ref{tab:kg_setting_yangstyle}。

\begin{table}[htbp]
    \centering
    \caption{KG-GCN模型主要实验参数}
    \label{tab:kg_setting_yangstyle}
    \begin{tabular}{ll}
        \toprule
        参数 & 取值 \\
        \midrule
        参数节点数量 & 24 \\
        残差窗口长度 & 60 \\
        GCN隐藏层数量 & 2 \\
        节点隐藏特征维数 & 32 \\
        Mask比例 & 15\% \\
        分类层隐藏维数 & 96 \\
        优化器 & AdamW \\
        权重衰减 & $1\times10^{-4}$ \\
        \bottomrule
    \end{tabular}
\end{table}

本文采用准确率（Accuracy）、精确率（Precision）、召回率（Recall）以及Macro-F1作为故障诊断评价指标，并利用混淆矩阵进一步分析不同故障类别的识别情况。准确率用于评价全部测试样本中被模型正确识别的比例，其计算公式为
\begin{equation}
Accuracy=\frac{\sum_{k=1}^{C}TP_k}{N_s},
\end{equation}
式中：$C$为故障类别数量；$N_s$为样本总数；$TP_k$表示第$k$类故障被正确识别的样本数量。

精确率用于评价模型预测为某一故障类别的样本中实际属于该类别的比例。第$k$类故障的精确率为
\begin{equation}
P_k=\frac{TP_k}{TP_k+FP_k},
\end{equation}
式中：$FP_k$表示其他类别被误判为第$k$类故障的样本数量。

召回率反映某一故障类别中被模型正确检出的样本比例，其计算公式为
\begin{equation}
R_k=\frac{TP_k}{TP_k+FN_k},
\end{equation}
式中：$FN_k$表示第$k$类故障被误判为其他类别的样本数量。

针对第$k$类故障，其F1值为精确率与召回率的调和平均：
\begin{equation}
F1_k=\frac{2P_kR_k}{P_k+R_k}.
\end{equation}
对各故障类别F1值进行等权平均得到Macro-F1：
\begin{equation}
F1_{macro}=\frac{1}{C}\sum_{k=1}^{C}F1_k.
\end{equation}
Macro-F1能够同时反映不同故障类别上的精确率和召回率，适合用于评价多类别故障诊断模型的综合识别性能。

\subsection{诊断结果与对比}
\label{subsec:results_yangstyle}

表~\ref{tab:main_result_yangstyle}列出了CNN、AE和KG-GCN三种模型的故障诊断结果。为与单次工程算例的结果呈现方式保持一致，表中给出固定随机种子123下的主实验结果；同时对三个随机种子进行重复训练，用于核验模型结果稳定性。从表中可以看出，KG-GCN模型准确率达到98.57\%，高于CNN的97.86\%和AE的95.54\%；在综合评价指标Macro-F1上，KG-GCN达到98.58\%，同样优于CNN的97.85\%和AE的95.54\%。此外，KG-GCN对一级抽汽支路泄漏和VHP汽封泄漏两类故障的召回率分别达到98.57\%和99.29\%。三次重复训练结果中，KG-GCN的平均准确率和Macro-F1分别为98.45\%和98.46\%，Macro-F1标准差仅为0.11\%，表明模型在不同随机初始化条件下仍保持较稳定的诊断性能。

\begin{table}[htbp]
    \centering
    \caption{CNN、AE和KG-GCN算法模型诊断结果对比}
    \label{tab:main_result_yangstyle}
    \resizebox{\textwidth}{!}{%
    \begin{tabular}{lccccccc}
        \toprule
        模型算法 & 准确率/\% & 精确率/\% & F0召回率/\% & F1召回率/\% & F2召回率/\% & F3召回率/\% & Macro-F1/\% \\
        \midrule
        CNN & 97.86 & 97.89 & 99.29 & 96.43 & 95.71 & 100.00 & 97.85 \\
        AE & 95.54 & 95.57 & 95.71 & 92.86 & 95.71 & 97.86 & 95.54 \\
        \textbf{KG-GCN} & \textbf{98.57} & \textbf{98.60} & \textbf{99.29} & \textbf{97.14} & \textbf{98.57} & \textbf{99.29} & \textbf{98.58} \\
        \bottomrule
    \end{tabular}}
\end{table}

KG-GCN模型将第三章知识图谱中的参数关联关系转化为图卷积计算路径，使参数特征不再仅按照矩阵中的局部位置进行提取，而是沿VHP本体、一级抽汽支路和一次再热接口等实际热力关系逐层传播。在具有多个共享参数的故障类别之间，这种结构约束能够减少与故障源关联较弱的局部噪声影响，并强化具有实际设备联系的多参数协同特征，因此在整体准确率和不同类别识别稳定性方面取得更好的结果。

为进一步分析不同模型的具体误诊情况，选取同一随机种子下的测试结果构造混淆矩阵。CNN、AE和KG-GCN对应的混淆矩阵分别为
\begin{equation}
C_{CNN}=\begin{bmatrix}
139&1&0&0\\
5&135&0&0\\
0&2&134&4\\
0&0&0&140
\end{bmatrix},
\end{equation}
\begin{equation}
C_{AE}=\begin{bmatrix}
134&1&1&4\\
6&130&4&0\\
3&3&134&0\\
2&1&0&137
\end{bmatrix},
\end{equation}
\begin{equation}
C_{KG-GCN}=\begin{bmatrix}
139&1&0&0\\
4&136&0&0\\
1&1&138&0\\
1&0&0&139
\end{bmatrix}.
\end{equation}
矩阵行表示真实类别，列表示模型预测类别。由三种模型的混淆矩阵可以看出，KG-GCN分类结果最集中于主对角线，整体误诊样本数量最少。其中，F0正常状态仅有1个样本被识别为F1，F2一级抽汽支路泄漏仅有2个样本发生误判，F3 VHP汽封泄漏仅有1个样本被误诊为F0。相比之下，CNN对F2的识别存在更明显的类别交叉，其中有4个F2样本被误诊为F3；AE的误诊分布更加分散，F0、F1和F2均出现多方向误判。综合对比可见，KG-GCN在多类别故障识别中具有更高的准确性和稳定性。

CNN模型主要依赖局部卷积窗口提取参数--时间特征。当不同故障在局部参数波动幅值或短时动态形态上具有较高相似性时，CNN难以显式利用参数之间的设备关联。例如，F2一级抽汽支路泄漏和F3 VHP汽封泄漏均可能在抽汽及VHP排汽区域产生状态变化，两类故障共享部分热力参数；当局部残差模式接近时，CNN容易产生类别交叉。AE模型的优化目标侧重于对输入状态进行压缩和重构，当不同故障在潜在空间内具有相近的重构特征时，模型可能将其映射至相似区域，因此在F0、F1和F2之间出现更多分散误判。

KG-GCN则利用知识图谱规定的参数连接路径进行特征传递。对于F2故障，一级抽汽口、逆止门前后以及1号高压加热器入口构成连续的抽汽支路关系；对于F1和F3故障，VHP本体压力、排汽状态及再热接口形成不同于抽汽支路的参数关联区域。图卷积在这些实际结构关系上逐层聚合残差特征，使模型能够同时识别异常幅值及其在系统中的空间位置和传播范围，从而减少仅依赖局部数值特征造成的误判。

考虑高质量故障标签在实际电站中通常较为稀缺，进一步对掩码预训练效果进行了补充验证。当仅使用20\%带标签故障窗口训练分类阶段时，直接从随机初始化训练KG-GCN的Macro-F1为94.48\%$\pm$0.93\%，采用无标签残差掩码预训练并固定图卷积编码器后，Macro-F1提高至96.67\%$\pm$0.27\%，平均提升2.19个百分点。这表明预训练能够利用未标注残差窗口提前学习参数之间的结构关联，在故障标签有限时为分类层提供更加稳定的特征表示。该结果与本文采用“无标签残差预训练--少量故障标签分类”的训练设计相一致。

\subsection{本章小结}

本章选取VHP作为代表性设备，对所提出的汽轮机系统KG-GCN故障诊断方法进行了算例验证。首先基于2024年6月真实健康运行残差和独立故障机理构造正常状态、VHP通流部件侵蚀、一级抽汽支路泄漏以及VHP汽封泄漏四类诊断样本，并按照自然日完成训练、验证和测试数据划分；随后采用准确率、精确率、召回率和Macro-F1对CNN、AE和KG-GCN进行对比。结果表明，在固定随机种子主实验中，KG-GCN准确率和Macro-F1分别达到98.57\%和98.58\%，高于CNN和AE，且混淆矩阵中的分类结果更加集中于主对角线。三次重复实验进一步表明，KG-GCN具有较好的结果稳定性；在仅使用20\%带标签故障样本时，掩码预训练可使Macro-F1提高2.19个百分点。上述结果验证了利用汽轮机领域知识关系约束多参数残差传播，并结合无标签残差预训练进行故障分类的可行性。需要指出的是，本章故障样本属于基于真实健康残差和独立机理构造的半物理场景，所得结论反映方法在该类构造场景下的有效性，后续仍需结合现场故障记录或经校准的热力仿真进一步验证工程泛化能力。
